\section{Parity of shear warping functions and diagonality of the averaged Jacobian}

\subsection*{Setup}
Let the beam cross-section be a bounded Lipschitz set \(\Omega\subset\mathbb{R}^2\) with coordinates \((x_2,x_3)\) and area \(A=|\Omega|\).
Assume \(\Omega\) is \emph{bisymmetric} with respect to the coordinate axes \(x_2=0\) and \(x_3=0\), which are principal axes of inertia:
\[
I_2=\int_\Omega x_2^2\,d\Omega,\qquad I_3=\int_\Omega x_3^2\,d\Omega,\qquad \int_\Omega x_2 x_3\,d\Omega=0.
\]

Let the vector of shear warping functions be
\[
\boldsymbol\phi=(\phi_2,\phi_3):\Omega\to\mathbb{R}^2,\qquad 
\nabla\boldsymbol\phi=
\begin{bmatrix}
\partial_{2}\phi_2 & \partial_{3}\phi_2\\
\partial_{2}\phi_3 & \partial_{3}\phi_3
\end{bmatrix}.
\]
Define the cross-section average \(\langle f\rangle := \frac{1}{A}\int_\Omega f\,d\Omega\).

We consider unit generalized shear along a principal axis (\(x_2\) or \(x_3\)). Material is linear elastic, isotropic with Young’s modulus \(E\), Poisson’s ratio \(\nu\), and shear modulus \(G=\tfrac{E}{2(1+\nu)}\).

% \subsection*{Classical Love/Saint-Venant construction}
For each shear direction \(\beta\in\{2,3\}\), introduce a harmonic potential \(\chi^{(\beta)}:\Omega\to\mathbb{R}\) solving the Love-type Neumann problem
\begin{equation*}
\left\{
\begin{aligned}
\Delta \chi^{(2)} &= 0 && \text{in }\Omega,\\
\frac{\partial \chi^{(2)}}{\partial n} &= -\,n_2\!\left(\tfrac{\nu}{2}x_2^2+\Big(1-\tfrac{\nu}{2}\Big)x_3^2\right) - n_3 (2+\nu)\,x_2x_3 && \text{on }\partial\Omega,\\
\int_\Omega \chi^{(2)}\,d\Omega &= 0, \\
\Delta \chi^{(3)} &= 0 && \text{in }\Omega,\\
\frac{\partial \chi^{(3)}}{\partial n} &= -\,n_3\!\left(\tfrac{\nu}{2}x_3^2+\Big(1-\tfrac{\nu}{2}\Big)x_2^2\right) - n_2 (2+\nu)\,x_2x_3 && \text{on }\partial\Omega,\\
\int_\Omega \chi^{(3)}\,d\Omega &= 0,
\end{aligned}\right.
\end{equation*}
where \(\boldsymbol n=(n_2,n_3)\) is the outward unit normal.

Define the shear-warping vector \(\boldsymbol\phi\) columnwise by
\begin{align}
\phi_2(x_2,x_3) &:= \frac{GA}{E I_2}\left[\chi^{(2)} + x_2 x_3^2
-\frac{1}{A}\int_\Omega\!\big(\chi^{(2)} + x_2 x_3^2\big)\,d\Omega
-\frac{x_2}{I_2}\int_\Omega\!x_2\big(\chi^{(2)} + x_2 x_3^2\big)\,d\Omega\right],\\
\phi_3(x_2,x_3) &:= \frac{GA}{E I_3}\left[\chi^{(3)} + x_3 x_2^2
-\frac{1}{A}\int_\Omega\!\big(\chi^{(3)} + x_3 x_2^2\big)\,d\Omega
-\frac{x_3}{I_3}\int_\Omega\!x_3\big(\chi^{(3)} + x_3 x_2^2\big)\,d\Omega\right].
\end{align}
These definitions enforce the gauges
\[
\int_\Omega \phi_\alpha\,d\Omega=0,\ 
\int_\Omega x_2\phi_\alpha\,d\Omega=0,\ 
\int_\Omega x_3\phi_\alpha\,d\Omega=0\ (\alpha=2,3).
\]

\subsection*{Reflection symmetries and parity}
Let the reflections \(\mathcal R_2(x_2,x_3)=(-x_2,x_3)\) and \(\mathcal R_3(x_2,x_3)=(x_2,-x_3)\).

\paragraph{Equivariance.}
The PDE \(\Delta\chi^{(\beta)}=0\), the Neumann boundary operator, and the gauges are linear and \emph{equivariant} under \(\mathcal R_2,\mathcal R_3\). Moreover, the prescribed Neumann data are polynomials in \((x_2,x_3)\) whose even/odd character is independent of \(\nu\). Therefore:
\begin{itemize}
\item For unit shear along \(x_2\): the Neumann data for \(\chi^{(2)}\) are \emph{odd} in \(x_2\) and \emph{even} in \(x_3\).
\item For unit shear along \(x_3\): the Neumann data for \(\chi^{(3)}\) are \emph{even} in \(x_2\) and \emph{odd} in \(x_3\).
\end{itemize}

\paragraph{Parity of \(\boldsymbol\phi\).}
Consider unit shear along \(x_2\). If \(\phi_2\) solves the gauged problem, then so do
\[
\tilde\phi_2(x):= -\,\phi_2(\mathcal R_2 x),\qquad \hat\phi_2(x):= \ \ \phi_2(\mathcal R_3 x),
\]
since the reflection flips/keeps the boundary data as required and the gauges are invariant. By uniqueness (modulo the imposed gauges), we conclude
\[
\phi_2(-x_2,x_3)=-\phi_2(x_2,x_3)\quad\text{(odd in \(x_2\))},\qquad
\phi_2(x_2,-x_3)=\ \ \phi_2(x_2,x_3)\quad\text{(even in \(x_3\))}.
\]
By the same argument, the companion warping \(\phi_3\) for the \(x_2\)-shear case satisfies
\[
\phi_3(-x_2,x_3)=\ \ \phi_3(x_2,x_3)\quad\text{(even in \(x_2\))},\qquad
\phi_3(x_2,-x_3)=-\phi_3(x_2,x_3)\quad\text{(odd in \(x_3\))}.
\]
For unit shear along \(x_3\), the roles of \(2\leftrightarrow3\) swap.

\paragraph{Independence of \(\nu\).}
The parity statements depend only on the even/odd structure of the polynomial Neumann data and symmetry of \(\Omega\); multiplying these polynomials by coefficients involving \(\nu\) does not change parity. Hence no restriction such as \(\nu=0\) is needed.

\subsection*{Diagonality of the averaged Jacobian \(\langle\nabla\boldsymbol\phi\rangle\)}
With the parities above (take the \(x_2\)-shear case without loss of generality):
\[
\partial_3\phi_2 \text{ is odd in }x_3 \ \Rightarrow\ \langle \partial_3\phi_2\rangle=\frac{1}{A}\int_\Omega \partial_3\phi_2\,d\Omega=0,
\]
\[
\partial_2\phi_3 \text{ is odd in }x_2 \ \Rightarrow\ \langle \partial_2\phi_3\rangle=\frac{1}{A}\int_\Omega \partial_2\phi_3\,d\Omega=0,
\]
\[
\partial_2\phi_2 \text{ is even in }x_2\text{ and }x_3 \ \Rightarrow\ \langle \partial_2\phi_2\rangle\ \text{may be nonzero},
\]
\[
\partial_3\phi_3 \text{ is even in }x_2\text{ and }x_3 \ \Rightarrow\ \langle \partial_3\phi_3\rangle\ \text{may be nonzero}.
\]
Therefore the averaged Jacobian is diagonal:
\[
\boxed{\ 
\Big\langle \nabla\boldsymbol\phi \Big\rangle
=
\begin{bmatrix}
\langle \partial_2\phi_2\rangle & 0\\[2pt]
0 & \langle \partial_3\phi_3\rangle
\end{bmatrix}.\
}
\]

\subsection*{Consequences}
For a bisymmetric section and principal-axis shear loading, supplying the full bidirectional warping data \(\boldsymbol\phi=(\phi_2,\phi_3)\) does \emph{not} induce spurious coupling under unidirectional loading: all off-diagonal averages vanish by parity, independent of the value of \(\nu\). Any constitutive mapping built from \(\langle \nabla\boldsymbol\phi\rangle\) (e.g., \((I+\langle\nabla\boldsymbol\phi\rangle)^{-1}\) or projection-based expressions) will therefore reduce to a diagonal form under these symmetries.
